% !TEX program = lualatex
\documentclass[10pt,a4paper]{article}
\usepackage[a4paper,margin=1.6cm,top=1.5cm,bottom=1.5cm]{geometry}
\usepackage{fontspec}
\usepackage{polyglossia}
\setmainlanguage{french}
\usepackage{colortbl}
\usepackage{booktabs,tabularx,array,xcolor,enumitem,amssymb,fancyhdr,eurosym,graphicx}
\setmainfont{Latin Modern Roman}
\setmonofont{Latin Modern Mono}[Scale=MatchLowercase]

\definecolor{accent}{RGB}{31,78,121}
\definecolor{pale}{RGB}{232,239,247}
\newcolumntype{C}[1]{>{\centering\arraybackslash}p{#1}}
\renewcommand{\arraystretch}{1.25}
\setlength{\parindent}{0pt}
\setlist{nosep,leftmargin=*}
\usepackage{titlesec}
\titleformat{\section}{\large\bfseries\color{accent}}{}{0pt}{}[\vspace{-2mm}\textcolor{accent}{\rule{\linewidth}{0.6pt}}]
\titlespacing*{\section}{0pt}{10pt}{5pt}

\pagestyle{fancy}
\fancyhf{}
\renewcommand{\headrulewidth}{0pt}
\fancyhead[L]{\small\textcolor{gray}{1\iere{} NSI -- Fiche de suivi enseignant}}
\fancyhead[R]{\small\textcolor{gray}{Projet Bandit Manchot}}
\fancyfoot[C]{\small\thepage}

\begin{document}

\begin{center}
{\LARGE\bfseries\color{accent} Fiche de suivi de projet}\\[2mm]
{\Large Projet : Bandit Manchot}
\end{center}

\medskip
\begin{tabularx}{\linewidth}{@{}l X@{}}
\textbf{Élèves} & 1. \dotfill\ 2. \dotfill\ 3. \dotfill \\[2mm]
\textbf{Classe / groupe} & \dotfill \hfill \textbf{Date de rendu} : \dotfill \\[2mm]
\textbf{Fichier(s)} & \texttt{Bandit\_Manchot\_Eleve.py} \\
\end{tabularx}

\section*{Suivi de l'avancement}
\begin{tabularx}{\linewidth}{@{}X l C{1.2cm} C{1.2cm} C{1.2cm} l@{}}
\toprule
\rowcolor{pale} Jalon & Tests unitaires & Commencé & Terminé & Tests OK & Date / remarque \\
\midrule
\texttt{choisir\_symbole} & doctest : 4 tests & $\square$ & $\square$ & $\square$ & \\
\texttt{fabriquer\_chaine} & doctest : 4 tests & $\square$ & $\square$ & $\square$ & \\
\texttt{compte\_symboles\_identiques} & doctest : 4 tests & $\square$ & $\square$ & $\square$ & \\
\texttt{presence\_symboles\_identiques\_multiples} & doctest : 5 tests & $\square$ & $\square$ & $\square$ & \\
\texttt{table\_gain} & doctest : 8 tests & $\square$ & $\square$ & $\square$ & \\
\texttt{saisir\_mise} &  & $\square$ & $\square$ & $\square$ & \\
\texttt{demander\_continuer} &  & $\square$ & $\square$ & $\square$ & \\
\texttt{afficher\_bandit} &  & $\square$ & $\square$ & $\square$ & \\
\texttt{main\_game} &  & $\square$ & $\square$ & $\square$ & \\
Partie facultative (High Score) &  & $\square$ & $\square$ & $\square$ & \\
\bottomrule
\end{tabularx}

\section*{Journal des séances}
\begin{tabularx}{\linewidth}{@{}C{2cm} X X@{}}
\toprule
\rowcolor{pale} Date & Travail observé / difficultés & Aides apportées, consignes données \\
\midrule
 & & \\[8mm]
\midrule
 & & \\[8mm]
\midrule
 & & \\[8mm]
\midrule
 & & \\[8mm]
\bottomrule
\end{tabularx}

\section*{Points de vigilance à contrôler}
\begin{itemize}
\item \texttt{get\_score} et \texttt{sauve\_score} n'ont pas été modifiées ; \texttt{HighScore.txt} est créé automatiquement s'il manque.
\item \texttt{saisir\_mise} : mise entre 10 et le pot ; refus d'un texte, d'un nombre négatif, d'un flottant, d'une mise supérieure au pot.
\item \texttt{table\_gain} : vérifier \texttt{777} ($\times 100$), \texttt{ΩΩΩ} ($\times 50$), un brelan de couleurs ($\times 20$), \texttt{Ω7Ω} ($\times 10$), \texttt{♠♠7} ($\times 5$), trois symboles différents ($\times 2$), sinon 0.
\item \texttt{main\_game} : départ à 500\,\euro, la partie s'arrête quand le pot est nul ou que le joueur abandonne ; la fonction renvoie le gain réel.
\item Information : avec cette table, l'espérance de gain est d'environ 2,59 fois la mise (très favorable au joueur).
\end{itemize}
\smallskip
Cas à tester : mise de 10, mise égale au pot, pot nul après une perte, plusieurs parties de suite, saisies absurdes.

\section*{Grille d'évaluation}
\begin{tabularx}{\linewidth}{@{}>{\raggedright\arraybackslash}p{4.6cm} X C{1.3cm} C{1.5cm} p{3.4cm}@{}}
\toprule
\rowcolor{pale} Critère & Détails & Barème & Obtenu & Observations \\
\midrule
\texttt{choisir\_symbole} & assertions efficaces & 1 & & \\
\texttt{fabriquer\_chaine} & assertions efficaces & 1 & & \\
\texttt{compte\_symboles\_identiques} & passage de tous les tests unitaires & 1 & & \\
\texttt{presence\_symboles\_identiques\_multiples} &  & 2 & & \\
\texttt{table\_gain} & passage de tous les tests unitaires & 1 & & \\
\texttt{saisir\_mise} & \emph{dumbproof} & 1 & & \\
\texttt{demander\_continuer} & \emph{dumbproof} & 1 & & \\
\texttt{afficher\_bandit} & affichage correct & 1 & & \\
\texttt{main\_game} & jeu a minima fonctionnel & 4 & & \\
Noms des variables clairs & pas de nom d'un seul caractère, sauf compteurs & 2 & & \\
Code commenté et clair & explications minimales & 2 & & \\
Améliorations, qualité du code &  & 1 & & \\
\midrule
\multicolumn{2}{l}{\textbf{Sous-total}} & \textbf{18} & & \\
\midrule
Partie facultative (High Score) & affichage correct et bonne saisie & 2 & & \\
\midrule
\multicolumn{2}{l}{\textbf{TOTAL}} & \textbf{20} & & \\
\bottomrule
\end{tabularx}


\section*{Bilan}
\begin{tabularx}{\linewidth}{@{}X@{}}
\textbf{Appréciation générale} :\par\medskip
\hrule\medskip\hrule\medskip\hrule\medskip
\end{tabularx}

\bigskip
\begin{tabularx}{\linewidth}{@{}X r@{}}
\textbf{Note du projet : \hspace{1.5cm} / 20} \quad $\Rightarrow$ \quad \textbf{Note sur 20 : \hspace{1.5cm}} & Date : \dotfill \\
\end{tabularx}

\end{document}
